% Part: first-order-logic
% Chapter: model-theory
% Section: isomorphism

\documentclass[../../../include/open-logic-section]{subfiles}

\begin{document}

\olfileid{mod}{bas}{iso}
\olsection{Isomorfe structuren}

!!{structure}s van de eerste orde kunnen op twee manieren op elkaar
lijken. Ten eerste kunnen in beide dezelfde !!{sentence}s waar zijn.
Zulke !!{structure}s noemen we \emph{elementair equivalent}. Structuren
kunnen echter sterk verschillen en toch dezelfde !!{sentence}s waar
maken: de ene kan bijvoorbeeld !!{enumerable} zijn en de andere niet.
Dat komt doordat veel eigenschappen van !!a{structure} niet in een taal
van de eerste orde kunnen worden uitgedrukt. Soms biedt de taal daarvoor
te weinig uitdrukkingsmiddelen; soms ligt de oorzaak bij fundamentele
beperkingen van de eerste-ordelogica, zoals de stelling van
L\"owenheim--Skolem. Er is daarom een tweede, strengere manier waarop
!!{structure}s op elkaar kunnen lijken: ze zijn fundamenteel hetzelfde
en verschillen alleen in de objecten waaruit ze bestaan, niet in hun
structurele kenmerken. Het begrip \emph{isomorfisme} maakt dit precies.

\begin{defn}
\ollabel{defn:elem-equiv} 
Gegeven twee !!{structure}s $\Struct{M}$ en $\Struct M'$ voor dezelfde
!!{language}~$\Lang{L}$ noemen we $\Struct{M}$ \emph{elementair
  equivalent met} $\Struct M'$, genoteerd als $\Struct{M} \equiv \Struct M'$,
dan en slechts dan als voor elke !!{sentence}~$!A$ van~$\Lang{L}$ geldt:
$\Sat{M}{!A}$ dan en slechts dan als $\Sat{M'}{!A}$.
\end{defn}

\begin{defn}
\ollabel{defn:isomorphism} Gegeven twee !!{structure}s $\Struct{M}$ en
$\Struct M'$ voor dezelfde !!{language}~$\Lang L$ noemen we
$\Struct{M}$ \emph{isomorf met}~$\Struct M'$, genoteerd als $\Struct{M}
\simeq \Struct M'$, dan en slechts dan als er een functie $h \colon
\Domain{M} \to \Domain{M'}$ bestaat waarvoor geldt:
\begin{enumerate}
\item $h$ is !!{injective}: als $h(x) =
  h(y)$, dan $x = y$; 
\item $h$ is !!{surjective}: voor elke $y \in \Domain{M'}$ bestaat
  er een $x \in \Domain{M}$ waarvoor $h(x) = y$;
\item \ollabel{defn:iso-const}voor elke !!{constant} $c$:
  $h(\Assign{c}{M}) = \Assign{c}{M'}$;
\item \ollabel{defn:iso-pred}voor elk $n$-plaatsig !!{predicate}~$P$:
  \[
  \tuple{a_1, \dots, a_n}\in \Assign{P}{M} \quad\text{iff}\quad
  \tuple{h(a_1), \dots, h(a_n)} \in \Assign{P}{M'};
  \]
\item \ollabel{defn:iso-func}voor elke $n$-plaatsige !!{function} $f$:
  \[
  h(\Assign{f}{M}(a_1, \dots, a_n)) =
  \Assign{f}{M'}(h(a_1), \dots, h(a_n)).
  \]
\end{enumerate}
\end{defn}

\begin{thm}
\ollabel{thm:isom}
Als $\Struct{M} \iso \Struct M'$, dan $\Struct{M} \elemequiv
\Struct{M'}$.
\end{thm}

\begin{proof}
Zij $h$ een isomorfisme van $\Struct{M}$ op $\Struct M'$. Voor elke
toekenning~$s$ is $h \circ s$ de samenstelling van $h$ en $s$, dus de
toekenning in $\Struct{M'}$ waarvoor $(h \circ s)(x) = h(s(x))$.
Met inductie naar $t$ en $!A$ kunnen de volgende sterkere beweringen
worden bewezen:
\begin{enumerate}
  \item[a.] $h(\Value{t}{M}[s]) = \Value{t}{M'}[h\circ s]$.
  \item[b.] $\Sat{M}{!A}[s]$ dan en slechts dan als $\Sat{M'}{!A}[h \circ s]$.
\end{enumerate}
De eerste bewering volgt met inductie naar de complexiteit van~$t$.
\begin{enumerate}
\item Als $t \ident c$, dan $\Value{c}{M}[s] = \Assign{c}{M}$ en
  $\Value{c}{M'}[h \circ s] = \Assign{c}{M'}$. Bijgevolg
  $h(\Value{t}{M}[s]) = h(\Assign{c}{M}) = \Assign{c}{M'}$ (volgens
  \olref{defn:iso-const} van \olref{defn:isomorphism}) $=
  \Value{t}{M'}[h \circ s]$.
\item Als $t \ident x$, dan $\Value{x}{M}[s] = s(x)$ en
  $\Value{x}{M'}[h \circ s] = h(s(x))$. Bijgevolg is $h(\Value{x}{M}[s]) =
  h(s(x)) = \Value{x}{M'}[h \circ s]$.
\item Als $t \ident f(t_1, \dots, t_n)$, dan
  \begin{align*}
    \Value{t}{M}[s] & = \Assign{f}{M}(\Value{t_1}{M}[s], \dots, \Value{t_n}{M}[s]) \quad\text{and}\\
  \Value{t}{M'}[h \circ s] & = \Assign{f}{M}(\Value{t_1}{M'}[h \circ
    s], \dots, \Value{t_n}{M'}[h \circ s]).
  \end{align*}
  De inductiehypothese luidt dat voor elke $i$, $h(\Value{t_i}{M}[s])
  = \Value{t_i}{M'}[h\circ s]$. Dus
  \begin{align}
    h(\Value{t}{M}[s]) 
    & = h(\Assign{f}{M}(\Value{t_1}{M}[s], \dots, \Value{t_n}{M}[s]) \notag\\
    & = \Assign{f}{M'}(h(\Value{t_1}{M}[s]), \dots,
    h(\Value{t_n}{M}[s])) \ollabel{iso-1}\\
    & = \Assign{f}{M'}(\Value{t_1}{M'}[h \circ s], \dots,
    \Value{t_n}{M'}[h \circ s]) \ollabel{iso-2}\\
    & = \Value{t}{M'}[h\circ s] \notag
  \end{align}
  Hierbij volgt \olref{iso-1} uit \olref{defn:iso-func} van
  \olref{defn:isomorphism} en \olref{iso-2} uit de inductiehypothese.
\end{enumerate}
Deel (b) wordt als oefening overgelaten.

Als $!A$ een zin is, zijn de toekenningen~$s$ en $h \circ s$
irrelevant en geldt $\Sat{M}{!A}$ dan en slechts dan als $\Sat{M'}{!A}$.
\end{proof}

\begin{prob}
Werk het bewijs van (b) van \olref[mod][bas][iso]{thm:isom} volledig
uit. Vermeld daarbij waar elk van de vijf eigenschappen wordt gebruikt
die de isomorfismen in \olref[mod][bas][iso]{defn:isomorphism}
karakteriseren.
\end{prob}

\begin{defn}
Een \emph{automorfisme} van een structuur $\Struct{M}$ is een isomorfisme
van $\Struct{M}$ op zichzelf.
\end{defn}

\begin{prob}
Toon voor elke structuur $\Struct{M}$ aan dat, als $X$ een definieerbare
deelverzameling van $\Struct{M}$ is en $h$ een automorfisme van
$\Struct{M}$, dan $X = \Setabs{h(x)}{x \in X}$ (dat wil zeggen: $X$ is
invariant onder $h$).
\end{prob}


\end{document}
